%! Sine and Cosine Function Values — Unit 3, Lesson 3.3 — Homework
\documentclass[10pt]{article}
\usepackage{precalculus-article}
\usepackage{precalculus-boxes}
\newcommand{\TallMath}[1]{$\displaystyle #1\rule[-1.4em]{0pt}{3.2em}$}

\begin{document}

\pageheader{Unit 3, Lesson 3.3}{Homework}
\namedateperiod

\vspace{0.1in}

\begin{notesbox}{Sine and Cosine Function Values --- Practice}

\textbf{1.} Complete the table with exact values. Use reference angles and quadrant sign conventions.

\vspace{0.06in}
\renewcommand{\arraystretch}{2.2}
\begin{tabularx}{\linewidth}{@{} >{\bfseries}l *{6}{>{\centering\arraybackslash}X} @{}}
\toprule
$\theta$ (rad) & $\dfrac{5\pi}{4}$ & $\dfrac{4\pi}{3}$ & $\dfrac{3\pi}{2}$ & $\dfrac{5\pi}{3}$ & $\dfrac{7\pi}{4}$ & $\dfrac{11\pi}{6}$ \\
\midrule
$\sin\theta$ & \underline{\hspace{1cm}} & \underline{\hspace{1cm}} & \underline{\hspace{1cm}} & \underline{\hspace{1cm}} & \underline{\hspace{1cm}} & \underline{\hspace{1cm}} \\
$\cos\theta$ & \underline{\hspace{1cm}} & \underline{\hspace{1cm}} & \underline{\hspace{1cm}} & \underline{\hspace{1cm}} & \underline{\hspace{1cm}} & \underline{\hspace{1cm}} \\
\bottomrule
\end{tabularx}
\renewcommand{\arraystretch}{1}

\vspace{0.14in}

\textbf{2.} A point $P$ lies on a circle of radius 8 centered at the origin, at angle
$\theta = \dfrac{7\pi}{6}$. Find the exact coordinates of $P$.

\vspace{0.8in}

$P = \Bigl($ \underline{\hspace{2.4cm}} $,$ \underline{\hspace{2.4cm}} $\Bigr)$

\vspace{0.14in}

\textbf{3.} Suppose $\sin\theta = -\dfrac{\sqrt{2}}{2}$ and $\cos\theta < 0$.
Identify $\theta$ in the interval $[0, 2\pi)$.

\vspace{0.6in}

$\theta = $ \underline{\hspace{2.2cm}}

\vspace{0.14in}

\textbf{4. AP-Style Multiple Choice.} Which expression gives the $y$-coordinate of
a point on a circle of radius $r$ at angle $\theta = \dfrac{5\pi}{3}$?

\medskip
\begin{tabular}{@{}lll@{}}
(A) $r \cdot \dfrac{\sqrt{3}}{2}$ \quad &
(B) $-r \cdot \dfrac{\sqrt{3}}{2}$ \quad &
(C) $r \cdot \dfrac{1}{2}$ \\[8pt]
(D) $-r \cdot \dfrac{1}{2}$ \quad &
(E) $r \cdot \dfrac{\sqrt{2}}{2}$ &
\end{tabular}

\vspace{0.1in}
Answer: \underline{\hspace{1.6cm}} \quad Justify: \underline{\hspace{6.5cm}}

\vspace{0.14in}

\textbf{5.} \textbf{Explain:} Using the definition $\cos\theta = x$ and $\sin\theta = y$
for the point $(x, y)$ on the unit circle at angle $\theta$, explain why $\cos(\pi + \theta) = -\cos\theta$
for any angle $\theta$. You may refer to the symmetry of the circle.

\vspace{1.1in}

\end{notesbox}

\begin{extensionbox}
\textbf{Extension (optional):} Using the same unit-circle reasoning as in Problem~5,
prove that $\sin(\pi - \theta) = \sin\theta$ for any angle $\theta$.
(\textit{Hint:} What happens to the $y$-coordinate when you reflect a point across the $y$-axis?)
\end{extensionbox}

\vspace{0.1in}

\begin{tcolorbox}[colback=goldbg, colframe=goldacc, arc=2mm,
    left=3mm, right=3mm, top=2mm, bottom=2mm]
\small\textbf{Preview of Lesson 3.4:} Now that we know exact $\sin\theta$ and $\cos\theta$ values
at key angles, we will use those values as ordered pairs to build and interpret the
\emph{graph} of $y = \sin x$ and $y = \cos x$. Bring your completed homework table --- those
values become the points on the graph!
\end{tcolorbox}

\end{document}
